\FifthArticle{5}
Aus der im vorhergehenden Artikel erörterten Reihenfolge der positiven und negativen Werte von $T$ und $U$, die bei jedem über $R$ hinausgehenden\SourceLeaf{086} Wert von $r$ gilt\HistoricalNote{14}, lässt sich nun folgern, dass innerhalb des Gebiets der kleineren Werte von $r$ gewisse Kreuzungen in diesen Anordnungen vorhanden sein müssen, die das Wesen unseres zu beweisenden Satzes enthalten. Ich werde die Beweisführung in einer der Geometrie der Lage entnommenen Form darstellen, weil sie dadurch die größte Anschaulichkeit und Einfachheit gewinnt. Im Grunde gehört aber der eigentliche Inhalt der ganzen Argumentation einem höheren, vom Räumlichen unabhängigen Gebiet der allgemeinen abstrakten Größenlehre an, dessen Gegenstand die durch Stetigkeit zusammenhängenden Größenkombinationen sind. Dieses Gebiet ist bisher noch wenig ausgebaut, und man kann sich darin auch nicht ohne eine von räumlichen Bildern entlehnte Sprache bewegen.

\FifthArticle{6}
Das ganze Gebiet der komplexen Größen wird durch eine unbegrenzte Ebene dargestellt, in der jeder Punkt mit den Koordinaten $t,u$ bezüglich zweier senkrecht aufeinander stehender Achsen als der komplexen Größe $x=t+iu$ entsprechend betrachtet wird. Bringt man diese komplexe Größe in die Form $x=r(\cos\rho+i\sin\rho)$, so bedeuten $r,\rho$ die Polarkoordinaten des entsprechenden Punktes. Die Gesamtheit aller komplexen Größen, für die $r$ denselben bestimmten Wert hat, wird demnach durch einen Kreis dargestellt, dessen Radius dieser Wert und dessen Mittelpunkt der Koordinatenursprung ist. Denjenigen dieser Kreise, für den $r$ um eine beliebig gewählte Differenz größer als $R$ ist, will ich mit $K$ bezeichnen, und mit $(1),\,(2),\,(3)\,.\,.\,.\,.\,(2n)$ diejenigen Punkte darauf, die den jeweils zwischen $\omega$ und $3\omega$, zwischen $5\omega$ und $7\omega$, zwischen $9\omega$ und $11\omega$ usw. bis zwischen $(8n-3)\omega$ und $(8n-1)\omega$ liegenden Werten von $\rho$ entsprechen, für die nach Artikel 4 $T=0$ wird. Man beachte dabei, dass für die Punkte $(1),\,(3),\,(5)$ usw. $U$ positiv, für die Punkte $(2),\,(4),\,(6)$ usw. hingegen negativ sein wird.
\SourceLeaf{087}
\Needspace{14\baselineskip}\FifthArticle{7}
Die Gesamtheit derjenigen Punkte in unserer Ebene, für die $T$ positiv ist, bildet zusammenhängende Flächenteile. Das wird schon deutlich, wenn man bedenkt, dass bei einem stetigen Übergang von einem Punkt zu einem anderen $T$ sich stetig ändert. Ebenso bilden sämtliche Punkte, für die $T$ negativ wird, zusammenhängende Flächenteile. Zwischen den Flächenteilen der ersten und denen der zweiten Art liegen Punkte, in denen $T=0$ wird. Nach der Natur der Funktion $T$ können diese Punkte nicht ebenfalls Flächenstücke, sondern nur Linien bilden, die auf der einen Seite die einen, auf der anderen die anderen Flächenteile begrenzen.

Der außerhalb von $K$ liegende Raum enthält $n$ Flächen der ersten Art, die mit ebenso vielen der zweiten Art abwechseln und von denen jede sich von einem Stück der Kreislinie $K$ aus zusammenhängend ins Unendliche erstreckt. Zugleich ist aber klar, dass jedes dieser Flächenstücke sich über die Kreislinie hinaus in den inneren Raum fortsetzt und dass für die weitere Gestalt folgende Fälle eintreten können.

1) Das betreffende, an einem Teil von $K$ beginnende Flächenstück endet für sich innerhalb der Kreisfläche. Seine äußere Begrenzung besteht dann nur aus zwei zusammenhängenden Stücken, von denen eines ein Bestandteil von $K$ ist und das andere innerhalb des Kreises liegt. In der beigefügten Figur, die sich auf eine Gleichung fünften Grades bezieht und in der die Vorzeichen von $T$ in die verschiedenen Flächenteile eingetragen sind, gehören drei der Flächen mit positivem $T$ zu diesem Fall. Die eine hat die Grenzlinien $10.1$ und $1.11.10$, die zweite $4.5$ und $5.12.4$, die dritte $6.7$ und $7.13.6$. Zwei ähnliche Flächenteile mit negativem $T$ sind vorhanden.

2) Das Flächenstück durchquert die Kreisfläche einfach so, dass es mit einem an einer anderen Stelle eintretenden eine zusammenhängende Fläche bildet. Die gesamte äußere Begrenzungslinie besteht dann aus vier Stücken, von denen zwei zur Kreislinie $K$ und die beiden anderen zum inneren Raum gehören. In unserer Figur tritt dieser Fall bei dem durch $2.3;$ $3.0.8;\ 8.9;\ 9.11.2$ begrenzten Flächenstück auf.

3) Das Flächenstück verzweigt sich im inneren Kreisraum einmal oder mehrmals so, dass es mit noch zwei oder mehr an anderen Stellen eintretenden eine zusammenhängende Fläche bildet. Deren gesamte äußere \SourceLeaf{088}Begrenzung besteht dann aus sechs, acht oder mehr Stücken in gerader Anzahl, die abwechselnd zur Kreislinie und zum inneren Raum gehören. In unserer Figur tritt dies bei einem Flächenteil ein, dessen Begrenzung durch die sechs Stücke $3.4;\ 4.12.5;\ 5.6;\ 6.13.7;\ 7.8;\ 8.0.3$ gebildet wird, in dem aber $T$ negativ ist.

\Needspace{19\baselineskip}\FifthArticle{8}

Bei einer vollständigen Aufzählung aller denkbaren Gestalten der in den inneren Kreisraum eintretenden Flächenteile müssten den angegebenen Fällen noch weitere Abwandlungen hinzugefügt werden. Wenn sich beispielsweise ein solcher Flächenteil in zwei Äste teilt, die sich im inneren Raum wieder vereinigen, könnte dieser Fall dem ersten oder zweiten Fall des vorhergehenden Artikels zugerechnet werden, je nachdem, ob die Fläche nach der Vereinigung im Inneren endet oder sich (ohne neue Teilung) bis zu einer anderen Stelle der Kreislinie fortsetzt. Die Gestalt der Fläche wäre dann nur durch den Einschluss einer nicht zu ihr gehörenden Insel verändert. Übrigens wäre es nicht schwer, streng zu beweisen, dass bei der besonderen Beschaffenheit der Funktion $T$ solche Abwandlungen \emph{überhaupt nicht möglich} sind. Für unseren Zweck ist das jedoch unnötig, da es nur auf die Folge der Stücke der \emph{äußeren} Begrenzung jedes der betreffenden Flächenteile ankommt (also derjenigen, in denen $T$ positiv ist).

Wir haben nämlich schon darauf hingewiesen, dass die Anzahl dieser Stücke immer gerade ist (zwei im ersten Fall des vorhergehenden Artikels, vier im zweiten, sechs oder mehr im dritten), wobei abwechselnd eines zur Kreislinie $K$ und eines zum inneren Raum gehört. Ferner ist klar: Wird jene äußere Begrenzungslinie stets in derselben Richtung durchlaufen, wofür hier diejenige gewählt werden soll, in der die Nummern der Punkte von $K$ zunehmen (also beispielsweise in unserer Figur so, dass die Fläche immer rechts von der Begrenzungslinie liegt), dann sind Anfangs- und Endpunkt eines zur Kreislinie gehörenden Stücks jeweils mit einer geraden und der um eins größeren ungeraden Zahl bezeichnet. Somit sind Anfangs- und Endpunkt jedes den inneren Raum durchlaufenden Stücks immer jeweils mit einer ungeraden und einer geraden Zahl bezeichnet.

Es steht also fest, dass von den $n$ Linien, die an einem mit einer ungeraden Zahl\SourceLeaf{089} bezeichneten Punkt von $K$ in den inneren Raum eintreten und in denen überall $T=0$ ist, jede diesen Raum auf eine ganz bestimmte Art\HistoricalNote{15} zusammenhängend durchläuft, bis sie an einer anderen, mit einer geraden Zahl bezeichneten Stelle wieder austritt. Da nun, wie schon oben (am Ende von Artikel 6) bemerkt, der Wert von $U$ an ihrem Anfangspunkt positiv und am Endpunkt negativ ist, muss wegen der Stetigkeit der Wertänderung notwendig an einem Zwischenpunkt $U=0$ werden. Dieser Punkt stellt dann eine Wurzel der Gleichung $X=0$ dar. Da die Anzahl solcher Linien $=n$ ist, ergeben sich auf diese Weise immer $n$ Wurzeln jener Gleichung.

\Needspace{18\baselineskip}\FifthArticle{9}

Wenn die genannten Linien den Kreisraum durchlaufen, ohne einander zu treffen, sind die so erhaltenen $n$ Wurzeln offenbar notwendig verschieden. Ein solches freies Durchlaufen findet sich in unserer Figur bei den Linien von $3$ nach $8$, von $5$ nach $4$ und von $7$ nach $6$; dazu gehören die durch die Punkte $0,\ 12,\ 13$ dargestellten Wurzeln. Haben hingegen zwei oder mehr solcher Linien einen Punkt gemeinsam, so ist dadurch zwar noch nicht notwendig, aber doch möglich, dass dieser Punkt zugleich derjenige ist, an dem $U=0$ wird. Dann fallen zwei oder mehr Wurzeln in eine zusammen, oder sie sind, wie man gewöhnlich sagt, einander gleich. In unserer Figur treffen die Linien $1.10$ und $9.2$ am Punkt $11$ zusammen, und dort wird zugleich $U=0$. Die Gleichung hat also außer den schon genannten drei verschiedenen noch zwei gleiche Wurzeln.

